\subsection{INVARIANTS OF A FINITE LINEAR GROUP: MODULAR PROPERTIES}

Let K be a commutative ring, V a K-module, and G a group acting on V. We know that every automorphism of V extends uniquely to an automorphism of the symmetric algebra \(S = \mathbf{S}(V)\) , and thus G acts on this algebra. If \(x \in S\) and \(g \in G\) , we denote by \(g_{S}.x\) the transform of x by g. Let R be the subalgebra \(S^{G}\) of S formed by the elements invariant under G.

Assume that G is finite, V is of finite type, and K is noetherian. Then S is an R-module of finite type, and R is a K-algebra of finite type (Commutative Algebra, Chap. V, §1, no. 9, Th. 2). Assume that S is integral and let N be its field of fractions. The field of fractions L of R is the set of elements of N invariant under G (loc. cit., Cor. of Prop. 23), so N is a Galois extension of L. Every element of N can be written z/t with \(z \in S\) and \(t \in R\) (loc. cit., Prop. 23). By Algebra, Chap. II, §7, no. 10, Cor. 3 of Prop. 26, the rank of the R-module S is {[}N : L{]}. Assume that G acts faithfully on V. The Galois group of N over L can then be identified with G, so {[}N : L{]} = Card G; thus

\[
\operatorname{rk} _ {\mathsf {R}} (\mathrm{S}) = [ \mathrm{N}: \mathrm{L} ] = \operatorname{Card} (\mathrm{G}).\tag{6}
\]

For any graded algebra \(A = A_{0} \oplus A_{1} \oplus \cdots \oplus A_{n} \oplus \cdots\) , denote by \(A_{+}\) the ideal \(\bigoplus_{n > 0} A_{n}\) .

THEOREM 1. Let K be a commutative field, V a finite dimensional vector space over K, S = S(V) the symmetric algebra of V, G a finite group of automorphisms of V, and R the graded subalgebra of S consisting of the elements invariant under G. Assume that G is generated by pseudo-reflections (§2, no. 1) and that q = Card(G) is coprime to the characteristic of K. Then the R-module S has a basis consisting of q homogeneous elements.

\begin{enumerate}
\def\labelenumi{\alph{enumi})}
\item
  Since every submodule of \(S/(R_{+}S)\) is free over \(R_0 = K\) , it is enough to show (in view of Algebra, Chap. II, §11, no. 4, Prop. 7) that the canonical homomorphism from \(R_{+} \otimes_{R} S\) to \(S\) is injective. For any R-module \(E\) , denote by \(T(E)\) the R-module \(\text{Ker}(R_{+} \otimes_{R} E \to E)\) (*in other words, \(T(E) = \text{Tor}_{1}^{R}(R/R_{+}, E)_*\) ). If \(E, E'\) are two R-modules and \(u\) is a homomorphism from \(E\) to \(E'\) , the homomorphism \(1 \otimes u\) from \(R_{+} \otimes E\) to \(R_{+} \otimes E'\) defines by restriction to \(T(E)\) a homomorphism from \(T(E)\) to \(T(E')\) that we denote by \(T(u)\) . If \(u'\) is a homomorphism from \(E'\) to an R-module \(E''\) , we have \(T(u' \circ u) = T(u') \circ T(u)\) . Thus, if G acts R-linearly on \(E\) , then G acts on \(T(E)\) .
\item
  The group G acts R-linearly on S, and hence also on T(S). Moreover, T(S) has a natural structure of graded S-module. We show first that, if \(g \in G\) , then g transforms every element x of T(S) into an element congruent to x modulo \(S_{1}T(S)\) . It is enough to do this when g is a pseudo-reflection. Then there exists a non-zero vector v in V such that \(g(x) - x \in Kv\) for all \(x \in V\) . Since V generates S, it follows that \(g_{S}\) acts trivially on S/Sv. Thus, for any \(y \in S\) , there exists an element \(h(y)\) in S such that
\end{enumerate}

\[
g _ {\mathrm{S}} (y) - y = h (y) v.
\]

Since S is integral and v is non-zero, this element is determined uniquely by y; it is immediate that h is an endomorphism of degree -1 of the R-module S. Thus, \(g_{S} - 1_{S} = m_{v} \circ h\) , where \(m_{v}\) denotes the homothety with ratio v in S. Hence,

\[
\mathrm{T} (g _ {\mathrm{S}}) - 1 _ {\mathrm{T(S)}} = \mathrm{T} (g _ {\mathrm{S}} - 1 _ {\mathrm{S}}) = \mathrm{T} (m _ {v}) \circ \mathrm{T} (h),
\]

the image of which is contained in \(vT(S)\) , proving our assertion.

\begin{enumerate}
\def\labelenumi{\alph{enumi})}
\setcounter{enumi}{2}
\tightlist
\item
  We show that any element of \(T(S)\) invariant under G is zero. Indeed, let Q be the endomorphism of the R-module S defined by
\end{enumerate}

\[
\mathrm{Q} (y) = q ^ {- 1} \sum_ {g \in G} g _ {\mathrm{S}} (y)
\]

for all \(y \in S\) . Then \(Q(S) = R\) . We can thus write \(Q = i \circ Q'\) , where \(Q'\) is a homomorphism from the R-module S to the R-module R and i denotes the canonical injection of R into S. Thus \(T(Q) = T(i) \circ T(Q')\) and \(T(Q') = 0\) since \(T(R) = \text{Ker}(R_{+} \otimes R \to R) = 0\) . Thus

\[
0 = \mathrm{T} (\mathrm{Q}) = q ^ {- 1} \sum_ {g \in \mathrm{G}} \mathrm{T} (g _ {\mathrm{S}}).
\]

But \(q^{-1} \sum_{g \in G} T(g_S)\) leaves fixed the elements of \(T(S)\) invariant under G. These elements are therefore zero.

\begin{enumerate}
\def\labelenumi{\alph{enumi})}
\setcounter{enumi}{3}
\tightlist
\item
  Assume that \(T(S) \neq 0\) . There exists in \(T(S)\) a homogeneous element \(u \neq 0\) of minimum degree. By b), u is invariant under G. By c), u = 0. This is a contradiction, so \(T(S) = 0\) . Q.E.D.
\end{enumerate}

Remarks. 1) It follows from Algebra, Chap. II, §11, no. 4, Prop. 7 that, if \((z_{1},\ldots ,z_{q})\) is a family of homogeneous elements of S whose canonical images in \(\mathrm{S / (R_+S)}\) form a basis of \(\mathrm{S / (R_+S)}\) over K, then \((z_{1},\ldots ,z_{q})\) is a basis of S over R.

\begin{enumerate}
\def\labelenumi{\arabic{enumi})}
\setcounter{enumi}{1}
\tightlist
\item
  Let \(g\) be a pseudo-reflection of \(V\) , whose order \(n \geqslant 2\) is finite and co-prime to the characteristic of \(K\) . By Maschke's theorem (Appendix, Prop. 2), \(V\) can be decomposed as \(D \oplus H\) , where \(H\) is the hyperplane consisting of the elements of \(V\) invariant under \(g\) and \(D\) is a line on which \(g\) acts by multiplication by a primitive \(n\) th root of unity. When \(K = R\) , this is possible only when \(n = 2\) , and \(g\) is then a reflection; in this case, the groups to which Th. 1 applies are the finite Coxeter groups. (For K = C, on the other hand, Th. 1 applies to certain groups that are not Coxeter groups.) \(^{3}\)
\end{enumerate}

THEOREM 2. Retain the assumptions and notation of Theorem 1.

\begin{enumerate}
\def\labelenumi{(\roman{enumi})}
\item
  There exists a graded vector subspace of \(S\) forming a complement to \(R_{+}S\) in \(S\) and stable under \(G\) .
\item
  Let U be such a complement. The canonical homomorphism from \(U \otimes_{K} R\) to S is an isomorphism of G-modules, and the representation of G in U (resp. S) is isomorphic to the regular representation of G on K (resp. R).
\end{enumerate}

Indeed, for any integer \(n \geqslant 0\) , the K-vector spaces \(S_n\) and \((R_+S) \cap S_n\) are stable under \(G\) , and it follows from Maschke's theorem (Appendix, Prop. 2) that there exists a G-stable complement \(U_n\) of \((R_+S) \cap S_n\) in \(S_n\) . Then \(\sum_{n \geqslant 0} U_n\) is a G-stable complement of \(R_+S\) in \(S\) , hence (i).

Let U be a graded vector subspace of S forming a complement of \(R_{+}S\) in S and stable under G. By Remark 1, every basis of the K-vector space U is also a basis of the R-module S, and consequently is a basis of the field of fractions N of S over the field of fractions L of R. Thus, the L-vector space N can be identified with \(U \otimes_{K} L\) . Since U is stable under G, this identification is compatible with the action of G. The group algebra \(L[G]\) of G over L can be identified with the algebra \(K[G] \otimes_{K} L\) . The Galois extension N of L admits a normal basis (Algebra, Chap. V, §10, no. 9, Th. 6), which can be interpreted as saying that N, considered as an \(L[G]\) -module, is isomorphic to the module of the regular representation of G over L. Since U is finite dimensional over K, it follows from the Appendix, Prop. 1, that the \(K[G]\) -module U is isomorphic to the regular representation of G over K. Our assertions follow immediately from this.

